%&pdflatex
\documentclass [12pt]{article}
\usepackage{authblk}
\usepackage{silence}
\usepackage{amsfonts}
\usepackage{amsmath}
\usepackage{wasysym}
\usepackage{hhline}
\usepackage{cite}
\usepackage{tocloft}
\usepackage[plainpages=false,pdfpagelabels]{hyperref}
\usepackage[nottoc]{tocbibind}

\WarningFilter{latexfont}{Font shape `U/wasy/b/n'}
\WarningFilter{latexfont}{Font shape `OMS/cmtt/m/n'}
\WarningFilter{latexfont}{Font shape `OMS/cmtt/bx/n'}

\setlength{\topmargin}{-2cm}
\setlength{\oddsidemargin}{0cm}
\setlength{\textheight}{23cm}
\setlength{\textwidth}{16cm}
\setlength{\parindent}{0pt}
\setlength{\parskip}{1.2pt}
\setlength{\cftparskip}{-2pt} % paragraph spacing across entries

\makeatletter
\newcommand{\tr}{\mathop{\mathrm{tr}}}
\makeatother

\title
    {
    Lattice Gauge Theory Symmetries
    \vspace*{-8pt}
    }
\date	{\small\today\vspace*{-30pt}}
\author	{Laurence G.~Yaffe}
\affil{%
    \small
    Department of Physics, University of Washington, Seattle WA 98195, USA
    \vspace*{-12pt}
    }%
\pagenumbering{roman}
\begin {document}
\maketitle
\tableofcontents
\setlength{\parskip}{5pt plus 2pt minus 1pt}
\pagenumbering{arabic}

\section{Introduction}

This note describes, in detail, the action of lattice gauge theory symmetries
acting on both the fundamental fields and on basic observables,
for theories, either Hamiltonian or Euclidean, defined on a simple cubic lattice.
Theories in Hamiltonian form are discussed first, followed by 
Euclidean theories.
This material expands on the presentation in Ref.~\cite{CSVAII}
(and corrects a few misprints).

Lattice Hamiltonians under consideration are the sum of
gauge and fermion pieces,
$
    H = H_{\rm gauge} + H_{\rm fermi}
$,
with the gauge Hamiltonian having the standard Kogut-Susskind form:
\begin{align}
    H_{\rm gauge}/N
    &=
    \tfrac 14 \lambda \sum_\ell \> \tr \, (E_\ell^2)
    + \lambda^{-1} \sum_p \> \tr\, (2 - U[\partial p] - U[\partial p]^\dagger) \,,
\label{eq:Hgauge}
\end{align}
where the sums run over all positively oriented links ($\ell$) and plaquettes ($p$).
The fermion discretization considered is the simplest staggered fermion formulation,
for which a single flavor fermion Hamiltonian takes the form:
\begin{align}
    H_{\rm fermi}
    &=
    \tfrac 12 \sum_{\langle xy \rangle} \>
    \psi^\dagger_x \, U_{\langle xy \rangle} \, \eta_{\langle xy \rangle} \, \psi_y
    + m \sum_x \> \sigma_x \, \psi^\dagger_x \, \psi_x \,,
\label{eq:Hfermi}
\end{align}
with the first sum running over all nearest-neighbor pairs of sites
(so that $\langle xy \rangle$ may be either a forward or backward directed link).
The ``staggering factor'' $\sigma_x$ appearing in the mass term is defined to equal
$+1$ on the even sublattice (i.e., sites displaced from the origin by an even
number of unit lattice translations), and $-1$ on the odd sublattice.%
\footnote
    {%
    The length of any compactified dimension must be even
    for the distinction between sublattices to be meaningful,
    and hence for the mass term to be well-defined.
    }

The ``flavor connection'' $\eta_{\langle xy \rangle}$ is a non-dynamical U$(1)$ connection
whose product around every plaquette equals minus one,
$\eta[\partial p] \equiv -1$;
this condition ensures recovery of the Dirac equation in the continuum limit.
This flavor connection may be chosen to be pure imaginary, and this is assumed
in the following.%
\footnote
    {%
    A standard form for the flavor connection is
    $
	\eta_{\langle xy \rangle} = i \prod_{\nu < \mu} (-1)^{x_\nu}
    $
    for a link $\langle xy \rangle$ running in the $+\mu$ direction
    starting at site $x$ with coordinates $\{ x_\mu \}$.
    But no such specific choice is ever needed.
    }
Fermion bilinears are most naturally defined with
the gauge connection and flavor connection appearing together,
so it is convenient to define the product
\begin{align}
    \hat U[\Gamma] \equiv U[\Gamma] \, \eta[\Gamma]
\label {eq:fermconnection}
\end{align}
as the fermion connection for any lattice path $\Gamma$.

\section  {Cubic lattice symmetries}\label{sec:Cubic}

Let $\Theta$ denote a cubic lattice symmetry
(some combination of discrete translations, rotations and/or reflections)
which maps sites to sites, $x \to \theta x$
and links to links, $\alpha \to \theta \alpha$.
The gauge fields (in a Hamiltonian formulation)
transform as:
\begin{align}
    (U_\alpha)_{ij} \xrightarrow{\Theta} (U_{\theta \alpha})_{ij} \,,
\qquad
    (E_\alpha)_{ij} \xrightarrow{\Theta} (E_{\theta \alpha})_{ij} \,.
\label{eq:thetagauge}
\end{align}
If a link $\alpha$ (or its transform) is backwards directed, so that $\alpha = \bar\ell$
for some forward directed link $\ell$, then
\begin{align}
    (U_{\bar\ell})_{ij} \equiv (U_\ell^\dagger)_{ij} = (U_\ell)_{ji}^* \,,
\qquad
    (E_{\bar\ell})_{ij} \equiv -(U_{\bar\ell} \, E_\ell \, U_\ell)_{ij} \,.
\end{align}
Due to the U(1) fermion flavor connection
appearing in the Hamiltonian (or action)
the transformation of fermion operators must include
a flavor connection running from the initial site to the transformed site,
\begin{align}
    \psi_x \xrightarrow{\Theta}
    \eta[\Gamma_{x,\theta x}] \> \psi_{\theta x} \,.
\label{eq:thetafermi}
\end{align}
The particular path $\Gamma_{x,\theta x}$ associated
with a lattice symmetry $\Theta$ is constructed by designating
(arbitrarily) some site $o$ as the origin,
choosing some path $\Gamma_{x,o}$ connecting site $x$ to the origin,
and choosing a path
$\Gamma_{o,\theta o}$ connecting the origin to its transform
under the lattice symmetry in question.
Then the complete flavor connection path is given by
\begin{equation}
    \Gamma_{x,\theta x} \equiv
    \Gamma_{x,o} \, \Gamma_{o,\theta o} \, (\overline{\theta\Gamma})_{\theta o,\theta x} \,,
\end{equation}
where $(\overline{\theta \Gamma})_{\theta o,\theta x}$ denotes the transformed, reversed
path which runs from $\theta o$ to $\theta x$.
This definition leaves the fermion hopping term invariant:
\begin{align}
    \tfrac 12 \sum_{\langle xy \rangle}
	\psi_x^\dagger \,
	\eta_{\langle xy \rangle} \,
	U_{\langle xy \rangle} \,
	\psi_y
    &\to
    \tfrac 12 \sum_{\langle \theta x \> \theta y \rangle}
	\psi_{\theta x}^\dagger \,
	\eta[\Gamma_{x,\theta x}]^\dagger \,
	\eta_{\langle xy \rangle} \,
	\eta[\Gamma_{y,\theta y}] \,
	U_{\langle \theta x \, \theta y \rangle} \,
	\psi_{\theta y}
    \nonumber\\ &=
    \tfrac 12 \sum_{\langle \theta x \> \theta y \rangle}
	\psi_{\theta x}^\dagger \,
	\eta_{\langle \theta x \> \theta y \rangle} \,
	U_{\langle \theta x \> \theta y \rangle} \, \psi_{\theta y}
    %\nonumber\\ &=
    =
    \tfrac 12 \sum_{\langle xy \rangle}
	\psi_x^\dagger \,
	\eta_{\langle xy \rangle} \,
	U_{\langle xy \rangle} \,
	\psi_y \,.
\end{align}

    To see this, note that
    \begin{align}
	    &\eta_{\langle \theta x \> \theta y \rangle}{}^\dagger \,
	    \eta[\Gamma_{x,\theta x}]^\dagger \,
	    \eta_{\langle xy \rangle} \,
	    \eta[\Gamma_{y,\theta y}]
	\nonumber\\ &\qquad {} =
	    \eta_{\langle \theta x \> \theta y \rangle}{}^\dagger \,
	    \eta[(\overline{\theta \Gamma})_{\theta o,\theta x}]^\dagger \,
	    \eta[\Gamma_{o,\theta o}]^\dagger \,
	    \eta[\Gamma_{x,o}]^\dagger \,
	    \eta_{\langle xy \rangle} \,
	    \eta[\Gamma_{y,o}] \,
	    \eta[\Gamma_{o,\theta o}] \,
	    \eta[(\overline{\theta \Gamma})_{\theta o,\theta y}]
	\nonumber\\ &\qquad {} =
	    \left(
	    \eta[(\overline{\theta \Gamma})_{\theta o,\theta x}] \,
	    \eta_{\langle \theta x \> \theta y \rangle} \,
	    \eta[(\overline{\theta \Gamma})_{\theta o,\theta y}]^\dagger
	    \right)^\dagger
	    \left(
	    \eta[\Gamma_{x,o}]^\dagger \,
	    \eta_{\langle xy \rangle} \,
	    \eta[\Gamma_{y,o}] \,
	    \right)
	\nonumber\\ & \qquad {} = 1 \,.
    \end{align}
    In the penultimate form, the second factor
    $
	    \left(
	    \eta[\Gamma_{x,o}]^\dagger \,
	    \eta_{\langle xy \rangle} \,
	    \eta[\Gamma_{y,o}]
	    \right)
    $
    is the flavor connection around a closed path running from the origin $o$ to $x$,
    then to $y$, and back to $o$,
    while the first factor
    $
	    \left(
	    \eta[(\overline{\theta \Gamma})_{\theta o,\theta x}] \,
	    \eta_{\langle \theta x \> \theta y \rangle} \,
	    \eta[(\overline{\theta \Gamma})_{\theta o,\theta y}]^\dagger
	    \right)^\dagger
    $
    is the inverse flavor connection of the symmetry transformed path running from
    $\theta o$ to $\theta x$, then to $\theta y$, and back to $\theta o$.
    These two factors cancel because the two closed paths are geometrically similar;
    the flavor connection around every plaquette gives an identical
    factor of $-1$, and any surface spanning the first loop will
    necessarily contain the same number of plaquettes as the symmetry
    transform of that surface spanning the second loop.

    Physical observables transform in the natural manner:%
    \begin{subequations}\label{eq:latttrans}
    \begin{align}
	\tr(U[\Gamma])
	&\xrightarrow{\Theta}
	\tr(U[\theta\Gamma]) \,,
    \\
	\tr({:}E_\alpha \, U[\Gamma]{:})
	&\xrightarrow{\Theta}
	\tr({:}E_{\theta\alpha} \, U[\theta\Gamma]{:}) \,,
    \\
	\tr({:}E_\alpha \, U[\Gamma] \, E_\beta \, U[\Gamma']{:})
	&\xrightarrow{\Theta}
	\tr({:}E_{\theta\alpha} \, U[\theta\Gamma] \, E_{\theta\beta} \, U[\theta\Gamma']{:}) \,,
    \\
	{:} \psi_x^\dagger \, \hat U[\Gamma_{xy}] \, \psi_y {:}
	&\xrightarrow{\Theta}
	{:} \psi_{\theta x}^\dagger \, \hat U[(\theta\Gamma)_{\theta x, \theta y}] \,
	    \psi_{\theta y} {:} \,,
	\label{eq:transferm}
    \\
	{:} \psi_x^\dagger \, \hat U[\Gamma] \, E_\alpha \, \hat U[\Gamma'] \psi_y {:}
	&\xrightarrow{\Theta}
	{:} \psi_{\theta x}^\dagger \, \hat U[\theta\Gamma] \,
	    E_{\theta\alpha} \, \hat U[\theta\Gamma'] \psi_{\theta y} {:} \,.
	\label{eq:transEferm}
    \end{align}
    \end{subequations}
    Here, ${:}\mathcal O{:}$ denotes the 
    normal ordered form of an observable containing
    electric field insertions and/or fermions.
    This ordering is defined to symmetrize the placement of electric field insertions
    at the beginning and end of the observable while also antisymmetrizing the ordering
    of fermions \cite{CSVAII}.  Specifically:
    \begin{subequations}
    \begin{align}
	:U[\Gamma]_{ij} (E_\alpha)_{jk} U[\Gamma']_{kl}:
	& \equiv
	\tfrac 12
	\bigl( (E_\alpha)_{jk} U[\Gamma]_{ij} U[\Gamma']_{kl}
    \nonumber
    \\ &\quad{}
	+
	U[\Gamma]_{ij} U[\Gamma']_{kl} (E_\alpha)_{jk} \bigr) ,
    \\[3pt]
	:U[\Gamma]_{ij} (E_\alpha)_{jk} U[\Gamma']_{kl} (E_\beta)_{lm} U[\Gamma'']_{mn}:
	& \equiv
	\tfrac 12
	\bigl(
	    (E_\alpha)_{jk} U[\Gamma]_{ij} U[\Gamma']_{kl} U[\Gamma'']_{mn} (E_\beta)_{lm}
    \nonumber
    \\ &\quad{}
	+
	    (E_\beta)_{lm} U[\Gamma]_{ij} U[\Gamma']_{kl} U[\Gamma'']_{mn} (E_\alpha)_{jk}
	\bigr) \,,
    \\
	:(\psi_x^\dagger)_i \, G[\Gamma]_{ij} (\psi_y)_j:
	& \equiv
	\tfrac 12
	\bigl(
	    (\psi_x^\dagger)_i (\psi_y)_j
	    -
	    (\psi_y)_j (\psi_x^\dagger)_i
	\bigr) \,
	{:} G[\Gamma]_{ij} {:}
    \nonumber
    \\ &=
	(\psi_x^\dagger)_i \, {:}G[\Gamma]_{ij}{:} \, (\psi_y)_j
	-\tfrac 12
	\delta_{xy} \> \tr ({:}G[\Gamma]{:}) \,,
    \label{eq:normordferm}
    \end{align}
    \end{subequations}
    where $G[\Gamma]$ denotes any string operator, possibly
    including electric field insertions.
    
    The simple form (\ref{eq:latttrans}) for the transformation for fermion bilinears
    is only possible because
    the connection $\hat U$ appearing in all fermionic operators is always
    understood to be the product (\ref{eq:fermconnection}) of the U($N$)
    gauge connection and the U(1) flavor connection.
    Verifying the transformation (\ref{eq:transferm})
    involves
    the same argument as above, with the flavor connection around the closed path
    $(\Gamma_{x,o})^\dagger \, \Gamma_{xy} \, \Gamma_{y,o}$
    canceling that around the
    transformed path
    $
	(\overline{\theta\Gamma})_{\theta o,\theta x} \,
	(\theta\Gamma)_{\theta x\, \theta y} \,
	(\overline{\theta\Gamma})_{\theta o,\theta y}^\dagger
    $,
    and likewise for (\ref{eq:transEferm}).

    A staggered sum over all lattice sites of a fermion bilinear
    transforms as
    \begin{align}
	\sum_x \,
	{:} \psi_x^\dagger \sigma_x \, \hat U[\Gamma_{xy}] \, \psi_y {:}
	&\xrightarrow{\Theta}
	\sigma_\theta
	\sum_{\theta x} \,
	{:} \psi_{\theta x}^\dagger \sigma_{\theta x} \,
	    \hat U[\theta\Gamma_{\theta x , \theta y}] \,
	    \psi_{\theta y} {:} \,,
    \end{align}
    with the overall factor of $\sigma_\theta \equiv \sigma_x \, \sigma_{\theta x}$
    equaling $-1$ for lattice symmetries which interchange even and odd sublattices,
    and $+1$ otherwise.
    The fermion mass term is such a staggered sum, and hence only translations by
    even numbers of lattice spacings
    are symmetries of the theory when massive fermions are present.

\section  {Charge conjugation}\label{sec:ChargeConj}

The action of charge conjugation on the gauge fields is given by
\begin{align}
    (U_\alpha)_{ij}
    \xrightarrow{C}
    (U_\alpha)_{ij}^*
    =
    (U_{\bar\alpha})_{ji} \,,\qquad
    (E_\alpha)_{ij}
    \xrightarrow{C}
    (E_\alpha)_{ij}^*
    =
    -(E_\alpha)_{ji} \,.
\label{eq:Cgauge}
\end{align}
The naive definition of the action of charge conjugation on fermions is
\begin{align}
    \psi_x
    \xrightarrow{C}
    \psi_x^\dagger \, c_x^* \,,
\label{eq:Cfermi}
\end{align}
with $c_x$ a charge conjugation phase which satisfies the
condition
\begin{align}
    c_x = - \eta_{\langle xy \rangle}^* \, c_y \, \eta_{\langle yx \rangle} \,.
\end{align}
With the assumed pure imaginary form for the flavor connection,
this charge conjugation phase may be set to unity.

The transformations (\ref{eq:Cgauge}) and (\ref{eq:Cfermi}) leave all terms
in the Hamiltonian invariant, except the fermion mass term which changes sign.
With multiple fermion flavors and a mass matrix which is unitarily
equivalent to its negative, one may define charge conjugation to include
a $\mathbb Z_2$ flavor exchange transformation and thereby produce a
symmetry of the massive theory.
But with a single staggered lattice fermion flavor, the only way to compensate
the change in sign of the fermion mass is to combine the naive charge
conjugation transformation with a lattice transformation which interchanges
even and odd sublattices, such as a translation $S$ by a single lattice spacing
(in some direction),
\begin{align}
    C \equiv C_{\rm naive} \, S \,.
\end{align}
This is the definition used.

When acting on physical observables, charge conjugation reverses the orientation
of the gauge connection, as well as producing a sign change for each
electric field insertion:
\begin{subequations}
\begin{align}
    \tr(U[\Gamma])
    &\xrightarrow{C}
    \tr(U[\Gamma])^*
    =
    \tr(U[\overline\Gamma]) \,,
\\[3pt]
    \tr({:} E_\alpha \, U[\Gamma]{:})
    &\xrightarrow{C}
    \tr({:} E_\alpha \, U[\Gamma]{:})^*
    =
    -\tr({:} E_\alpha \, U[\overline\Gamma]{:}) \,,
\\[3pt]
    \tr({:}E_\alpha \, U[\Gamma] \, E_\beta \, U[\Gamma']{:})
    &\xrightarrow{C}
    \tr({:}E_\alpha \, U[\Gamma] \, E_\beta \, U[\Gamma']{:})^*
    =
    \tr({:}E_\alpha \, U[\overline\Gamma'] \, E_\beta \, U[\overline\Gamma]{:}) \,,
\end{align}
\end{subequations}
where $\overline\Gamma$ denotes the reverse of a path $\Gamma$.
Fermion bilinears summed over all lattice sites (i.e., momentum zero) transform as:
\begin{subequations}\label{eq:Cfermobs}
\begin{align}
    \sum_x \,
    {:} \psi_x^\dagger \, \hat U[\Gamma_{xy}] \, \psi_y {:}
    &\xrightarrow{C}
    -(-1)^{|\Gamma|}
    \sum_y \,
    {:} \psi_y^\dagger \, \hat U[\overline\Gamma_{yx}] \,
	\psi_x {:} \,,
    \label{eq:Cferm}
\\
    \sum_x \,
    {:} \psi_x^\dagger \, \hat U[\Gamma] \, E_\alpha \, \hat U[\Gamma'] \psi_y {:}
    &\xrightarrow{C}
    (-1)^{|\Gamma|}
    \sum_y \,
    {:} \psi_y^\dagger \, \hat U[\overline\Gamma'] \,
	E_\alpha \, \hat U[\overline\Gamma] \psi_x {:} \,.
    \label{eq:CEferm}
\end{align}
\end{subequations}
Note that charge conjugation preserves the normal ordered forms of observables.
The factor of $-(-1)^{|\Gamma|}$ in (\ref{eq:Cferm})
arises from the reversal of the implicit flavor connection which
(with a pure imaginary flavor connection) gives a $-1$ for each link
in the path $\Gamma$, plus an additional minus sign from the
anticommutation of the fermion fields, while
(\ref{eq:CEferm}) has one further minus sign from the electric field
transformation.
%
For staggered (i.e., momentum $\pi$) fermion bilinears, the relations
analogous to (\ref{eq:Cferm}) and (\ref{eq:CEferm}) read:
\begin{subequations}\label{eq:Cfermstag}
\begin{align}
    \sum_x \,
    {:} \psi_x^\dagger \, \sigma_x \,\hat U[\Gamma_{xy}] \, \psi_y {:}
    &\xrightarrow{C}
    (-1)^{|\Gamma|}
    \sum_x \,
    {:} \psi_y^\dagger \, \hat U[\overline\Gamma_{yx}] \,
	\sigma_x \, \psi_x {:}
\nonumber \\ & \;\;\quad {}=
    \sum_y \,
    {:} \psi_y^\dagger \, \sigma_y \,\hat U[\overline\Gamma_{yx}]  \,
	\psi_x {:} \,,
\\
    \sum_x \,
    {:} \psi_x^\dagger \, \sigma_x \,\hat U[\Gamma] \, E_\alpha \, \hat U[\Gamma'] \psi_y {:}
    &\xrightarrow{C}
    -\sum_y
    {:} \psi_y^\dagger \, \sigma_y \,\hat U[\overline\Gamma'] \,
	E_\alpha \, \hat U[\overline\Gamma] \psi_x {:} \,,
\end{align}
\end{subequations}
where the change by a factor of $-(-1)^{|\Gamma|}$ in these expressions
relative to (\ref{eq:Cfermobs}) comes from  rewriting
$\sigma_x = (-1)^{|\Gamma|} \, \sigma_y$, together with the
additional minus sign from the unit shift in the definition of $C$.

\section {Time reversal}\label{sec:TimeReversal}

Time reversal (in the Hamiltonian formulation)
is an anti-unitary transformation which acts on the gauge fields as:
\begin{subequations}
\begin{align}
    (U_\alpha)_{ij}
    \xrightarrow{T}
    (U_{\bar\alpha})_{ji} \,,\quad
    (E_\alpha)_{ij}
    \xrightarrow{T}
    -(E_\alpha)_{ji} \,.
\label{eq:Tgauge}
\end{align}
while the naive definition of time reversal on fermions is
\begin{align}
    (\psi_x)_i
    \xrightarrow{T}
    (\psi_x^\dagger)_i \,\sigma_x \,.
\label{eq:Tfermi}
\end{align}
\end{subequations}
This transformation leaves
the fermion kinetic term invariant,
\begin{align}
    \sum_{\langle xy\rangle} \>
    \psi_x^\dagger \, \eta_{\langle xy \rangle} \,  U_{\langle xy\rangle} \, \psi_y
    &\xrightarrow{T}
    \sum_{\langle xy\rangle} \>
    \psi_x \, \sigma_x \,
    \eta_{\langle xy \rangle}^* \,
    U_{\langle xy\rangle}^* \,
    \sigma_y \, \psi_y^\dagger
    =
    \sum_{\langle yx\rangle} \>
    \psi_y^\dagger \, \eta_{\langle yx \rangle} \, U_{\langle yx\rangle} \, \psi_x \,.
\end{align}
The product of staggering factors cancels the minus sign
arising from anticommuting the fermions, while the
anti-unitary nature of the transformation causes the non-dynamical
flavor connection in this (and any other fermion bilinear)
to be complex-conjugated, which is equivalent to reversing its direction.
The pure gauge terms in the Hamiltonian are likewise invariant.

The transformation (\ref{eq:Tfermi}) acts on
the fermion mass term as
\begin{align}
    m \sum_x
    \psi_x^\dagger \, \sigma_x \, \psi_x
    &\xrightarrow{T}
    -m^* \sum_x \>
    \psi_x^\dagger \, \sigma_x \, \psi_x
    =
    -m \sum_x \>
    \psi_x^\dagger \, \sigma_x \, \psi_x \,.
\end{align}
Since the fermion mass must be real (for a Hermitian Hamiltonian),
the sign of the mass term flips.
Therefore, just as for complex conjugation, in the presence of
fermions a time reversal transformation will be combined with a single
lattice spacing translation,
\begin{align}
    T \equiv T_{\rm naive} \, S \,,
\end{align}
to produce a valid symmetry.

When acting on our set of physical observables, the action
of time reversal is identical to the effect of charge conjugation:
\begin{subequations}\label{eq:Ttrans}
\begin{align}
    \tr(U[\Gamma])
    &\xrightarrow{T}
    \tr(U[\overline\Gamma]) \,,
\\[3pt]
    \tr({:}E_\alpha \, U[\Gamma]{:})
    &\xrightarrow{T}
    -\tr({:}E_\alpha \, U[\overline\Gamma]{:}) \,,
\\[3pt]
    \tr({:}E_\alpha \, U[\Gamma] \, E_\beta \, U[\Gamma']{:})
    &\xrightarrow{T}
    \tr({:}E_\alpha \, U[\overline\Gamma'] \, E_\beta \, U[\overline\Gamma]{:}) \,.
\\[5pt]
    \sum_x \,
    {:} \psi_x^\dagger \, \hat U[\Gamma_{xy}] \, \psi_y {:}
    &\xrightarrow{T}
    -(-1)^{|\Gamma|}
    \sum_x \,
    {:} \psi_y^\dagger \, \hat U[\overline\Gamma_{yx}] \,
	\psi_x {:} \,,
\\
    \sum_x \,
    {:} \psi_x^\dagger \, \sigma_x \, \hat U[\Gamma_{xy}] \, \psi_y {:}
    &\xrightarrow{T}
    \sum_y \,
    {:} \psi_y^\dagger \, \sigma_y \, \hat U[\overline\Gamma_{yx}] \,
	\psi_x {:} \,,
\\
    \sum_x \,
    {:} \psi_x^\dagger \, \hat U[\Gamma] \, E_\alpha \, \hat U[\Gamma'] \psi_y {:}
    &\xrightarrow{T}
    (-1)^{|\Gamma|}
    \sum_y \,
    {:} \psi_y^\dagger \, \hat U[\overline\Gamma'] \,
	E_\alpha \, \hat U[\overline\Gamma] \psi_x {:} \,,
\\
    \sum_x \,
    {:} \psi_x^\dagger \, \sigma_x \, \hat U[\Gamma] \, E_\alpha \, \hat U[\Gamma'] \psi_y {:}
    &\xrightarrow{T}
    -\sum_y
    {:} \psi_y^\dagger \, \sigma_y \, \hat U[\overline\Gamma'] \,
	E_\alpha \, \hat U[\overline\Gamma] \psi_x {:} \,.
\end{align}
\end{subequations}

\section {Hermitian conjugation}\label{sec:HermitianConj}

While not a symmetry per se, it is useful to compare the above transformations with
the action of Hermitian conjugation on observables:
\begin{subequations}\label{eq:Hermconj}
\begin{align}
    \tr(U[\Gamma])^\dagger
    &\xrightarrow{\dagger}
    \tr(U[\overline\Gamma]) \,,
\\[3pt]
    \tr({:}E_\alpha \, U[\Gamma]{:})
    &\xrightarrow{\dagger}
    \tr({:}E_\alpha \, U[\overline\Gamma]{:}) \,,
\\[3pt]
    \tr({:}E_\alpha \, U[\Gamma] \, E_\beta \, U[\Gamma']{:})
    &\xrightarrow{\dagger}
    \tr({:}E_\alpha \, U[\overline\Gamma'] \, E_\beta \, U[\overline\Gamma]{:}) \,.
\\[5pt]
    \sum_x \,
    {:} \psi_x^\dagger \, \hat U[\Gamma_{xy}] \, \psi_y {:}
    &\xrightarrow{\dagger}
    \sum_y \,
    {:} \psi_y^\dagger \, \hat U[\overline\Gamma_{yx}] \,
	\psi_x {:} \,,
\\
    \sum_x \,
    {:} \psi_x^\dagger \sigma_x \, \hat U[\Gamma_{xy}] \, \psi_y {:}
    &\xrightarrow{\dagger}
    (-1)^{|\Gamma|}
    \sum_y \,
    {:} \psi_y^\dagger \sigma_y \, \hat U[\overline\Gamma_{yx}] \,
	\psi_x {:} \,,
\label{eq:fermstagdag}
\\
    \sum_x \,
    {:} \psi_x^\dagger \, \hat U[\Gamma] \, E_\alpha \, \hat U[\Gamma'] \psi_y {:}
    &\xrightarrow{\dagger}
    \sum_y \,
    {:} \psi_y^\dagger \, \hat U[\overline\Gamma'] \,
	E_\alpha \, \hat U[\overline\Gamma] \psi_x {:} \,,
\\
    \sum_x \,
    {:} \psi_x^\dagger \sigma_x \, \hat U[\Gamma] \, E_\alpha \, \hat U[\Gamma'] \psi_y {:}
    &\xrightarrow{\dagger}
    (-1)^{|\Gamma|}
    \sum_y
    {:} \psi_y^\dagger \sigma_y \, \hat U[\overline\Gamma'] \,
	E_\alpha \, \hat U[\overline\Gamma] \psi_x {:} \,.
\label{eq:Efermstagdag}
\end{align}
\end{subequations}
The factors of $(-1)^{|\Gamma|}$ in (\ref{eq:fermstagdag})
and (\ref{eq:Efermstagdag}) result from moving the staggering factor
to the other end of the bilinear.

These Hermitian conjugation transformations differ from the corresponding
actions of time reversal, or charge conjugation, by at most a minus sign.
Consequently, the expectation value of any observable $\mathcal O$ for which
charge conjugation equals either plus or minus Hermitian conjugation,
$
    C^\dagger \, \mathcal O C = \pm \mathcal O^\dagger
$,
in any state $|u\rangle$ which is charge conjugation invariant
must be either purely real or purely imaginary, since
\begin{align}
    \langle u | \mathcal O |u \rangle^*
    =
    \langle u | \mathcal O^\dagger |u \rangle
    =
    \pm \langle u | C^\dagger \mathcal O \, C |u \rangle
    =
    \pm \langle u | \mathcal O |u \rangle \,.
\label{eq:Creality}
\end{align}
Comparing (\ref{eq:Ttrans}) to (\ref{eq:Hermconj}) one sees that
observables with real expectation values (in any $C$-invariant state)
are:
\begin{itemize}\advance\itemsep -6pt
\item
    Wilson loops;
\item
    Wilson loops with two electric field insertions;
\item
    Odd-length, unstaggered fermion bilinears;
\item
    Even-length, staggered fermion bilinears;
\item
    Even-length, unstaggered fermion bilinears with one electric field insertion;%
    \footnote
	{%
	``Length'' means the number of links in the bilinear;
	the electric field insertion does not count.
	}
\item
    Odd-length, staggered fermion bilinears with one electric field insertion.
\end{itemize}
while observables with imaginary expectation values are:
\begin{itemize}\advance\itemsep -6pt
\item
    Wilson loops with one electric field insertion;
\item
    Even-length, unstaggered fermion bilinears;
\item
    Odd-length, staggered fermion bilinears;
\item
    Odd-length, unstaggered fermion bilinears with one electric field insertion;
\item
    Even-length, staggered fermion bilinears with one electric field insertion.
\end{itemize}

In Hamiltonian theories, the relation between time reversal and Hermitian conjugation
implies that in any time-reversal invariant state
the expectation values of observables $\mathcal O$ for which
$
    T \mathcal O T^{-1} = - \mathcal O^\dagger
$
(which are the same as observables for which
$
    C^\dagger \mathcal O C = - \mathcal O^\dagger
$)
actually vanishes in any $T$-invariant state.
This follows from combining this relation with
the defining relation for anti-unitary transformations,
$
    \langle T \Psi | T \chi \rangle
    =
    \langle \Psi | \chi \rangle^*
$,
and the invariance of the state
$
    |u\rangle = T |u \rangle
$,
to obtain:
\begin{align}
    \langle u | \mathcal O | u \rangle
    =
    \langle T u | \mathcal O | T u \rangle
    =
    \langle T u | \mathcal O^\dagger | T u \rangle^*
    =
    - \langle T u | T \mathcal O T^{-1} | T u \rangle^*
    =
    - \langle u | \mathcal O | u \rangle \,.
\end{align}
The first equality is $T$-invariance of $|u\rangle$,
the second is the definition of Hermitian conjugation,
the third is the assumed property of $\mathcal O$,
and the last is the defining relation of time-reversal
with $\langle \Psi| = \langle u|$ and
$|\chi\rangle = \mathcal O T^{-1}|T u\rangle = \mathcal O |u\rangle$.
Consequently, in any $C$- and $T$-invariant state, observables for which
charge conjugation would imply an imaginary expectation value actually have
vanishing expectation values.

\section {Implications}\label{sec:Consequences}

The above symmetry relations have numerous implications regarding the structure
of the large-$N$ phase space and associated dynamics.
These include:
\begin{itemize}\advance\itemsep -6pt
\item
    Time reversal invariance implies
    that the $N = \infty$ dynamics has a natural configuration space
    corresponding to the subspace of time-reversal invariant points within the
    $N = \infty$ phase space, with these points arising from the large-$N$ limit
    of (equivalence classes of) $T$-invariant coherent states.

\item
    Any function on phase space may be decomposed into components which
    are either even or odd under time reversal.
    In particular, phase space coordinates may be partitioned into those
    which are even under time reversal, i.e., configuration space coordinates,
    or those which are odd, i.e.., classical momenta.

\item
    Classical coordinates correspond to expectation values of the operators
    with real expectation values (in $C$-invariant states),
    while classical momenta correspond to expectation values of the
    operators for which charge conjugation implies imaginary expectation values,
    but whose expectation values actually vanish on the classical configuration space
    due to time reversal invariance.

\item
    Elements of the coherence algebra, which generate deformations in phase space,
    may be written as a linear combination of basis generators,
    $\mathfrak g = \sum_a c^a \, \hat e_a$, with the coefficients $\{ c^a \}$
    serving as Riemann normal coordinates in the neighborhood of any point
    in phase space.
    The generators $\hat e_a$
    may be partitioned into those which are time-reversal
    even, and those which are time-reversal odd.
    Time-reversal even generators
    are anti-Hermitian linear combinations of
    loops with a single electric field insertion,
    even-length unstaggered fermion bilinears,
    and odd-length staggered fermion bilinears;
    the coefficients $\{ c^a \}$ of these generators
    serve as configuration space coordinates.
    Correspondingly,
    time-reversal odd generators
    are anti-Hermitian linear combinations of
    Wilson loops,
    odd-length unstaggered fermion bilinears,
    and even-length staggered fermion bilinears;
    the coefficients $\{ c^a \}$ of these generators
    serve as classical momenta.

\item
    In order to be properly anti-Hermitian, coherence group generators composed
    of linear combinations of operators for which
    $C^\dagger \mathcal O C = \mathcal O^\dagger$,
    namely Wilson loops, odd-length unstaggered bilinears,
    and even-length staggered bilinears,
    must have imaginary coefficients when projected onto $C$-even
    symmetry representations, and real coefficient when projected onto
    $C$-odd representations.
    Because time reversal is anti-unitary, both types of such generators are
    time-reversal odd.

    The opposite situation applies to generators composed
    of linear combinations of operators for which
    $C^\dagger \mathcal O C = -\mathcal O^\dagger$,
    namely Wilson loops with an electric field insertion,
    even-length unstaggered bilinears,
    and odd-length staggered bilinears.
    Such generators must have real coefficients when projected
    onto $C$-even symmetry representations, and imaginary
    coefficients when projected onto $C$-odd representations.
    Both types of such generators are time-reversal even.

\end{itemize}

These various categories of coherence group generators are
shown schematically in Table \ref{tab:gens}, classified according
to their behavior under $C$ and $T$ transformations.
Of course, each such category can be further projected onto
particular irreducible representations of the lattice cubic
symmetry point group.

\begin{table}
\begin{center}
\setlength{\tabcolsep}{4pt}
\hspace*{-10pt}%
\begin{tabular}{ccccc}
$T$ & $C$ & gauge generators & \multicolumn{2}{c}{fermion generators}
\\[3pt]
+ & + & \phantom i $\tr (E \, U[\Gamma] - E \, U[\overline\Gamma])$
	& $ (\psi^\dagger \, \sigma\, \hat U[\Gamma_{\rm o}]\, \psi +
	      \psi^\dagger \, \sigma\, \hat U[\overline\Gamma_{\rm o}]\, \psi) $
	& $ (\psi^\dagger \, \hat U[\Gamma_{\rm e}]\, \psi -
	      \psi^\dagger \, \hat U[\overline\Gamma_{\rm e}]\, \psi) $

\\[3pt]
$-$ & + & $i \tr ( U[\Gamma] + U[\overline\Gamma])$
	& $ i(\psi^\dagger \, \hat U[\Gamma_{\rm o}]\, \psi +
	      \psi^\dagger \, \hat U[\overline\Gamma_{\rm o}]\, \psi) $
	& $ i(\psi^\dagger \, \sigma\, \hat U[\Gamma_{\rm e}]\, \psi +
	      \psi^\dagger \, \sigma\, \hat U[\overline\Gamma_{\rm e}]\, \psi )$

\\[3pt]
+ & $-$ & $i \tr (E \, U[\Gamma] + E \, U[\overline\Gamma])$
	& $ i(\psi^\dagger \, \sigma\, \hat U[\Gamma_{\rm o}]\, \psi -
	      \psi^\dagger \, \sigma\, \hat U[\overline\Gamma_{\rm o}]\, \psi )$
	& $ i(\psi^\dagger \, \hat U[\Gamma_{\rm e}]\, \psi +
	      \psi^\dagger \, \hat U[\overline\Gamma_{\rm e}]\, \psi )$

\\[3pt]
$-$ & $-$ & $ \phantom i \tr (U[\Gamma] - U[\overline\Gamma])$
	& $ (\psi^\dagger \, \hat U[\Gamma_{\rm o}]\, \psi -
	      \psi^\dagger \, \hat U[\overline\Gamma_{\rm o}]\, \psi )$
	& $ (\psi^\dagger \, \sigma\, \hat U[\Gamma_{\rm e}]\, \psi -
	      \psi^\dagger \, \sigma\, \hat U[\overline\Gamma_{\rm e}]\, \psi )$
\end{tabular}%
\hspace*{-10pt}
\end{center}
\caption
    {%
    Schematic forms of anti-Hermitian coherence group generators,
    classified according to their behavior under charge conjugation and
    time reversal symmetries.
    $\Gamma_e$ denotes an even-length lattice path and
    $\Gamma_o$ an odd-length path,
    overbars indicate path reversal,
    and $\sigma$ is the lattice staggering factor at the location of the preceding
    conjugate fermion field.
    Coefficients of the listed time-reversal even generators serve as Riemann normal
    coordinates within the classical configuration space of the fermion large-$N$ phase
    space, while coefficients of time-reversal odd generators serve as conjugate momenta.
    Only the first category of fully symmetric generators need be used for minimization.
    \label{tab:gens}
    }
\end{table}


\section {Euclidean theory symmetries}\label{sec:Euclidean}

Lattice actions for
Euclidean gauge theories under consideration 
are the sum of gauge and fermion contributions,
$A = A_{\rm gauge} + A_{\rm fermi}$, with the gauge action
having a standard Wilson form,
\begin{align}
    A_{\rm gauge}/N
    =
    \lambda^{-1} \sum_p \> \tr (2 - U[\partial p] - U[\partial p]^\dagger) \,.
\label{eq:Agauge}
\end{align}
Obviously, this gauge action is real provided the gauge coupling $\lambda$ is
real, which is assumed.

The most economical discretization for a single (massive) fermion flavor 
has just $N$ pairs (for U($N$) fundamentals) 
of independent Grassmann generators,
$\{ \chi{}_x^{\,i},\, \overline\chi{}_x^{\,i}\}$,
residing on each lattice site.
The fermion lattice action may be defined in a manner
completely analogous to the Hamiltonian form (\ref{eq:Hfermi}),
\begin{align}
    A_{\rm fermi}
    =
    \tfrac 12 \sum_{\langle xy \rangle} \>
    \overline\chi_x \, U_{\langle xy \rangle} \, \eta_{\langle xy \rangle} \, \chi_y
    + m \sum_x \> \sigma_x \, \overline\chi_x \, \chi_x \,,
\label{eq:Afermiold}
\end{align}
where, once again, the U(1) flavor connection $\eta_{\langle xy\rangle}$ is defined
to be pure imaginary and to have a product of minus one around every plaquette,
$\eta[\partial p] = -1$.
More abstractly, this fermion action is a quadratic form,
$
    A_{\rm fermi} = (\overline\chi, D \,\chi)
$,
where $D$ is a discretized Dirac operator with non-vanishing matrix elements
only for coincident or nearest-neighbor sites:
\begin{align}
    (D)_{xy}^{ij} \equiv
    \tfrac 12 \, U_{\langle xy \rangle}^{ij} \, \eta_{\langle xy \rangle}
    + m \, \sigma_x \, \delta_{xy} \, \delta^{ij} \,.
\end{align}
Because both the U($N$) gauge connection and the flavor connection 
are unitary connections
(i.e., $U_{\langle xy \rangle}^\dagger = U_{\langle yx \rangle}$
and $\eta_{\langle xy \rangle}^* = \eta_{\langle yx \rangle}$)
this Dirac operator is Hermitian,
provided the fermion mass $m$ is real, which is assumed.
Hermiticity of $D$ ensures reality of the fermionic partition function,
\begin{align}
    Z_{\rm fermi} \equiv \int d\overline \chi \, d\chi \> e^{-A_{\rm fermi}}
    = \det (D) \,.
\end{align}

This may be equivalently expressed as invariance of
the action $A_{\rm fermi}$ under a complex conjugation transformation
(an involutive automorphism)
which performs ordinary complex conjugation on all
commuting numbers while simultaneously interchanging Grassmann variables,
\begin{align}
    \chi_x \xrightarrow{*}  \overline\chi_x \,,\quad
    \overline\chi_x \xrightarrow{*}  \chi_x \,,
\label{eq:fermirealityold}
\end{align}
and flipping the order of all products of Grassmann variables.
%
Any statistical ensemble,
such as the equilibrium ensemble
$
    \rho_{\rm eq} \equiv Z^{-1} \, e^{-A}
$
or those generated by the
Euclidean large-$N$ coherence group acting on an initial uncorrelated ensemble,
must satisfy this same reality condition,
\begin{align}
    \rho(\{U_\ell\},\{\overline\chi_x\},\{\chi_x\})^*
    =
    \rho(\{U_\ell\},\{\overline\chi_x\},\{\chi_x\}) \,.
\label{eq:rhoreal}
\end{align}
The relevant physical observables in Euclidean theories are just
Wilson loops and fermion bilinears.
The reality condition (\ref{eq:rhoreal})
implies that expectation values of these observables satisfy
the basic Hermiticity relations:
\begin{align}
    \langle \tr(U[\Gamma]) \rangle^*
    =
    \langle \tr(U[\overline\Gamma]) \rangle \,,
    \qquad
    \big\langle \overline \chi_x \, \hat U[\Gamma_{xy}] \, \chi_y \big\rangle^*
    =
    \big\langle \overline \chi_y \, \hat U[\overline\Gamma_{yx}] \, \chi_x \big\rangle \,,
\label{eq:bilinearsold}
\end{align}
which are identical in form to the corresponding relations in Hamiltonian theories.

Ref.~\cite{CSVAII} chose to discuss Euclidean theories with
a fermion action of the form (\ref{eq:Afermiold}),
as this choice makes the action of symmetry transformations
as similar as possible between Hamiltonian and Euclidean formulations
and also makes the reality condition (\ref{eq:fermirealityold})
and its consequences (\ref{eq:bilinearsold})
as simple as possible.
However, the presence of a lattice staggering factor in the mass term
is an inconvenience which, in particular, obscures the fact that
the Euclidean theory (\ref{eq:Afermiold}) actually possess full lattice
translation invariance, including translations by an odd number of lattice spacings
(unlike the Hamiltonian theory).

The independence of the Grassmann generators
$\overline\chi$ and $\chi$ in the Euclidean formulation
allows one to make a change of variables involving
independent phase rotations of these variables, and this
enables one to make a field redefinition which will 
simplify the form of lattice symmetry transformations.
Let
\begin{align}
    \epsilon_x \equiv
    \begin{cases}
	1 \,, & x \mbox { even}; \\
	i \,, & x \mbox{ odd}\,,
    \end{cases}
\end{align}
denote a square root of the staggering factor $\sigma_x$,
and define
\begin{align}
    \overline\chi_x \equiv \overline\psi_x \, \epsilon_x^* \,,
    \qquad
    \chi_x \equiv i \sigma_x \, \epsilon_x \, \psi_x \,.
\label{eq:redef}
\end{align}
This field redefinition converts the lattice staggering factor in the
fermion mass term into a factor of $i$, while leaving the hopping term unchanged,
so the fermion action becomes
\begin{align}
    \widetilde A_{\rm fermi}
    =
    \tfrac 12 \sum_{\langle xy \rangle} \>
    \overline\psi_x \, U_{\langle xy \rangle} \,
    \eta_{\langle xy \rangle} \, \psi_y
    + i \, m \sum_x \> \overline\psi_x \, \psi_x \,.
\label{eq:Afermi}
\end{align}
The key benefit of this form for the fermion action
is that it obviates the
need for separate representations of staggered and
unstaggered versions of fermion bilinears.
The fermion action (\ref{eq:Afermi}) agrees with the form used
in most other literature on Euclidean staggered lattice fermions
(see, e.g., Ref.~\cite{Kronfeld:2007ek}),
except for an overall factor of $i$ in the action (\ref{eq:Afermi})
(hidden in the imaginary $\eta_{\langle xy \rangle}$ in
the hopping term and explicit in the mass term).
The presence, or absence, of this overall 
factor of $i$ is physically irrelevant, but its
inclusion makes $\eta_{\langle xy \rangle}$ a conventional
unitary Abelian connection which is convenient for the following symmetry analysis.

The fermion action (\ref{eq:Afermi}) remains a quadratic form,
$
    \widetilde A_{\rm fermi} = (\overline\psi,\widetilde D\, \psi)
$,
but with a transformed Dirac operator
$\widetilde D \equiv -i \epsilon^* D \epsilon$
now obeying a Hermiticity condition of the form
$\widetilde D^\dagger = - \sigma \, \widetilde D \, \sigma$.
This Hermiticity condition may also be expressed as invariance of
the action (\ref{eq:Afermi}) under a complex conjugation transformation
in which the redefined Grassmann fields transform as
\begin{align}
    \psi_x \xrightarrow{*}  i \, \overline\psi_x \, \sigma_x \,,\quad
    \overline\psi_x \xrightarrow{*}  i \, \sigma_x \, \psi_x \,,
\label{eq:fermireality}
\end{align}
(while also inverting the order of Grassmann products).
In terms of the redefined Grassmann fields,
the reality condition (\ref{eq:bilinearsold}) on fermion bilinear
expectation values becomes
\begin{align}
    \big\langle \overline \psi_x \, \hat U[\Gamma_{xy}] \, \psi_y \big\rangle^*
    &=
    -(-1)^{|\Gamma|}
    \big\langle \overline \psi_y \, \hat U[\overline\Gamma_{yx}] \, \psi_x \big\rangle \,.
\label{eq:fermirealityobs}
\end{align}

Before proceeding to an explicit discussion of the implications
of different symmetries,
recall that coherence group generators for Euclidean theories
contain first order derivatives acting on the fundamental fields.
These are defined as
\begin{align}
    E_{\ell,ij} \equiv
    \frac 1N \, U_{\ell,ik} \, \frac{\partial}{\partial U_{\ell,jk}} \,,
    \quad
    \overline \upsilon_{x,i} \equiv \frac{\partial}{\partial \psi_{x,i}} \,,
    \quad
    \upsilon_{x,i} \equiv \frac{\partial}{\partial \overline\psi_{x,i}} \,,
\end{align}
with the Grassmann derivatives $\upsilon_{x,i}$ and $\overline\upsilon_{x,i}$
satisfying the basic anticommutation relations
$
    \{ \overline\upsilon_{x,i},\, \psi_{y,j} \} = \delta_{xy} \, \delta_{ij}
$
and
$
    \{ \upsilon_{x,i},\, \overline\psi_{y,j} \} = \delta_{xy} \, \delta_{ij}
$,
with other anticommutators vanishing.
Invariance of these relations under a
complex conjugation transformation (\ref{eq:fermireality})
implies that the Grassmann derivatives behave as
\begin{align}
    \upsilon_x \xrightarrow{*}  -i \, \overline\upsilon_x \, \sigma_x \,,\quad
    \overline\upsilon_x \xrightarrow{*}  -i \, \sigma_x \, \upsilon_x 
\label{eq:fermirealityv}
\end{align}
under complex conjugation.

Cubic lattice symmetry transformations act identically to
the Hamiltonian formulation transformations
(\ref{eq:thetagauge}) and (\ref{eq:thetafermi}),
with the understanding that link derivatives $\{ E_\ell \}$ transform in
the same manner as electric field operators of the Hamiltonian theory
(and hence are denoted identically),
while the conjugate fermion $\overline\psi_x$ transforms as does 
the Hermitian conjugate $\psi^\dagger_x$ of the Hamiltonian theory.%
\footnote
    {%
    Time-reversal, if one so designates some lattice direction, is just
    part of the lattice symmetry group.
    }
The Grassmann derivatives $\overline \upsilon_x$ and $\upsilon_x$
transform in the same manner as do $\overline \psi_x$ and $\psi_x$, respectively.

Since the fermion mass term no longer contains a lattice staggering factor,
both gauge (\ref{eq:Agauge}))
and fermion (\ref{eq:Afermi}) actions in the Euclidean formulation
are invariant under odd lattice translations as well as even.%
\footnote
    {
    Note, however, that to define a flavor connection satisfying $\eta[\partial p] = -1$
    on every plaquette, it is necessary that the length of any compactified
    lattice direction be even.
    }

The action of charge conjugation on gauge fields is identical to the
Hamiltonian form (\ref{eq:Cgauge}), while the action on fermion fields
and their derivatives
may be defined as
\begin{equation}
    \psi_x \xrightarrow{C} \overline\psi_x \,\sigma_x \,,\quad
    \overline\psi_x \xrightarrow{C} - \sigma_x \, \psi_x \,,\quad
    \upsilon_x \xrightarrow{C} -\overline \upsilon_x \,\sigma_x \,,\quad
    \overline \upsilon_x \xrightarrow{C} \sigma_x \, \upsilon_x \,.
\label{eq:Cfermieuclid}
\end{equation}

In the mass term, the explicit minus sign from the transformation cancels
the minus sign from anticommuting the fermion fields back to their original order.
In the hopping term, the factor of $-\sigma_x \sigma_y$ equals one  since $x$ and $y$
are nearest-neighbor sites, while the minus sign from anticommuting the fermion
fields cancels the minus sign from reversing the flavor connection.
The transformations for the Grassmann derivatives $\upsilon_x$ and $\overline\upsilon_x$
follow from their definition as derivatives.

Acting on Wilson loops or fermion bilinears,
charge conjugation merely reverses the paths:
\begin{align}
    \tr(U[\Gamma])
    \xrightarrow{C}
    \tr(U[\overline\Gamma]) \,,
    \qquad
    \overline\psi_x \, \hat U[\Gamma_{xy}] \, \psi_y
    \xrightarrow{C}
    \overline\psi_y \, \hat U[\overline\Gamma_{yx}] \, \psi_x \,.
\label{eq:Ctransobs}
\end{align}
Just as for the hopping term of the action, the product of 
lattice staggering factors from the transformation (\ref{eq:Cfermieuclid})
cancels the minus signs from reversal of the flavor connection.
Charge conjugation acts on
fermion bilinear coherence group generators as
\begin{align}
    \overline\upsilon_x \, \hat U[\Gamma_{xy}] \, \psi_y
    \xrightarrow{C}
    -\overline\psi_y \, \hat U[\overline\Gamma_{yx}] \, \upsilon_x \,,
    \qquad
    \overline\psi_x \, \hat U[\Gamma_{xy}] \, \upsilon_y
    \xrightarrow{C}
    -\overline\upsilon_y \, \hat U[\overline\Gamma_{yx}] \, \psi_x \,.
\label{eq:Ctransgens}
\end{align}

Combining the charge conjugation relations (\ref{eq:Ctransobs}) with
the complex conjugation relations
(\ref{eq:bilinearsold}) and (\ref{eq:fermirealityobs})
shows that in any $C$-invariant state:
\begin{itemize}\advance\itemsep -3pt
\item $\langle \tr(U[\Gamma]) \rangle$ is real;
\item $\langle \overline\psi \,\hat U[\Gamma_e] \,\psi\rangle$ is imaginary;
\item $\langle \overline\psi \,\hat U[\Gamma_o] \,\psi\rangle$ is real;
\end{itemize}
with $\Gamma_e$ and $\Gamma_o$ denoting arbitrary paths of even or odd length,
respectively.
So the ``price'' for making the field redefinition (\ref{eq:redef})
is that all even length bilinears now have pure imaginary expectation values,
while odd length bilinears have real expectations.
The compensating benefit of using the form (\ref{eq:Afermi}) for the
fermion action is that physical observables now have very simple
behavior under all lattice symmetry transformations as well as
charge conjugation.

Table \ref{tab:euclidgens} shows the schematic forms of
coherence group generators for Euclidean gauge theories
when projected onto charge conjugation eigenstates.
The overall factors of $i$ in odd-length $C$-even
and even-length $C$-odd generators are
dictated by the requirement that deformations of the
statistical ensemble generated by these generators,
when multiplied by real coefficients, preserve
the reality condition for the statistical measure.

\begin{table}[t]
\begin{center}
\begin{tabular}{ccccc}
$C$ & gauge generators & \multicolumn{2}{c}{fermion generators}
\\[3pt]
  $+$ & $\tr\, (E U[\Gamma] - E U[\overline\Gamma])$
      & $(\overline \upsilon \, \hat U[\Gamma_e] \, \psi -
	\overline \psi \, \hat U[\overline \Gamma_e] \, \upsilon)$
      & $i\,(\overline \upsilon \, \hat U[\Gamma_o] \, \psi -
	\overline \psi \, \hat U[\overline \Gamma_o] \, \upsilon)$
\\[3pt]
  $-$ & $i \, \tr\, (E U[\Gamma] + E U[\overline\Gamma])$
      & $i\,(\overline \upsilon \, \hat U[\Gamma_e] \, \psi +
	\overline \psi \, \hat U[\overline \Gamma_e] \, \upsilon)$
      & $(\overline \upsilon \, \hat U[\Gamma_o] \, \psi +
	\overline \psi \, \hat U[\overline \Gamma_o] \, \upsilon)$
\end{tabular}
\end{center}
\caption
    {%
    Schematic forms of coherence group generators for Euclidean gauge theories,
    classified according to their behavior under charge conjugation.
    $\Gamma_e$ denotes an even-length lattice path and
    $\Gamma_o$ an odd-length path, while overbars indicate path reversal.
    \label{tab:euclidgens}
    }
\end{table}

\begin{thebibliography}{99}

%\cite{CSVAII}
\bibitem{CSVAII}
T.~A.~Dickens, U.~J.~Lindqwister, W.~R.~Somsky and L.~G.~Yaffe,
``The coherent state variational algorithm. 2. Implementation and testing,''
\href{https://doi.org/10.1016/0550-3213(88)90234-9}
{Nucl.~Phys.~B \textbf{309}, 1 (1988)}.
%doi:10.1016/0550-3213(88)90234-9
%20 citations counted in INSPIRE as of 01 Dec 2024

%\cite{Kronfeld:2007ek}
\bibitem{Kronfeld:2007ek}
A.~S.~Kronfeld,
``Lattice gauge theory with staggered fermions: how, where, and why (not),''
\href{https://doi:10.22323/1.042.0016}
{PoS \textbf{LATTICE2007}, 016 (2007)}
[arXiv:0711.0699 [hep-lat]].
%114 citations counted in INSPIRE as of 23 Sep 2025

\end{thebibliography}
\end{document}
